\documentclass[11pt,a4paper]{article}

\usepackage{amsmath}
\usepackage{amsfonts}
\usepackage{amssymb}
\usepackage{booktabs}
\usepackage{graphicx}
\usepackage{pgfplots}
\usepackage{tikz}
\usepackage{url}

\pgfplotsset{
    unbounded coords=discard,
    filter discard warning=false,
    compat=1.18
}

% Activate high-performance PGFPlots caching framework
\usepgfplotslibrary{external}
\tikzexternalize[prefix=tikz-cache/]

% Custom conditional plotting filter definition macro
\pgfplotsset{
    discard if not/.style 2 args={
        filter point/.code={
            \edef\val{\thisrow{#1}}
            \edef\checkval{#2}
            \ifx\val\checkval
            \else
                \pgfmapsetdiscardpoint
            \fi
        }
    }
}

\title{\textbf{SpectralBL Multi-Campaign Analysis:}\\
       \large Cross-Ecosystem Stable Boundary Layer Comparison}
\author{\textbf{Automated Analytics Pipeline Engine}\\
        \small SpectralBL-Analytics Framework}
\date{\today}

\begin{document}

\maketitle

\begin{abstract}
This report provides a comparative, low-rank structural analysis of stable boundary layer (SBL) regimes across multiple field campaigns and ecosystems. By projecting high-dimensional physical profile measurements from diverse sources (CASES-99, GABLS3, AMERIFLUX, SMEAR, NEON, ICOS, SHEBA) onto a unified low-rank continuous phase space, we reconstruct the underlying attractor manifold and analyze how different environments and forcings shape boundary layer organization. This framework provides an information-theoretic approach to diagnosing boundary layer complexity and mapping regime transitions across heterogeneous datasets.
\end{abstract}

\section{Introduction and Physical Context}

Traditional boundary layer parameterizations rely heavily on local gradient profiles and standard Monin-Obukhov Similarity Theory (MOST). However, in the very stable boundary layer (vSBL), the coupling between the surface layer and the upper air column is frequently severed by the development of a strong thermal inversion. The atmosphere enters a regime dominated by radiative cooling, low-level jet (LLJ) development, global intermittency, and submesoscale internal gravity waves.

In these conditions, local turbulence metrics such as $\mathrm{Ri}_g$ become unreliable diagnostics. To track structural reconfigurations across diverse ecosystems and measurement networks, the SpectralBL pipeline extracts continuous spatial pattern modes of the boundary layer state. The resulting low-rank coordinates $(\eta_1, \eta_2, \eta_3)$ remain dynamically active throughout stable and super-critical conditions, providing comparative insights into how different environments organize their stability structures.

\newpage

% Injection hooks for Mustache-rendered sub-sections
% templates/attractor.tex.mustache
% AUTO-GENERATED Attractor Report Section
% This file is rendered by Mustache template from trajectory data
% DO NOT EDIT MANUALLY

\section{Attractor Dynamics and Low-Dimensional Phase Space}

\subsection{Overview}

The boundary layer is analyzed as a low-dimensional dynamical system through projection onto dominant singular vectors. The following section presents the reconstructed phase space trajectory for campaign \texttt{ ALL } using the v0.1 baseline.

\begin{description}
    \item[Campaign:] ALL
    \item[Baseline Version:] v0.1
    \item[Total Samples:] 6682
    \item[Mean Entropy:] $H_{\text{mean}} = 0.223$
\end{description}

\subsection{3D Phase Space Projection}

Figure~\ref{fig:attractor_ALL} visualizes the three-dimensional trajectory $(\eta_1, \eta_2, \eta_3)$. To highlight the true multi-regime orbital evolution, timesteps are connected sequentially by a continuous historical path line. The trajectory is colored using a high-contrast colormap driven by the \textbf{Attractor Spin} diagnostic ($\Omega(t)$), explicitly isolating the directionality of regime transitions.

\begin{figure}[htbp]
    \centering
    \begin{tikzpicture}
        \begin{axis}[
            width=0.9\textwidth,
            height=0.55\textwidth,
            view={45}{24},
            grid=major,
            xlabel={$\eta_1$},
            ylabel={$\eta_2$},
            zlabel={$\eta_3$},
            title={Low-Rank Attractor Trajectory (\texttt{ ALL })},
            colormap/jet,
            colorbar,
            point meta=explicit,
            colorbar style={title={$\Omega(t)$}, title style={yshift=5pt}},
            tick label style={font=\small},
            label style={font=\small},
            title style={font=\small}
        ]
            % 1. DRAW THE ORBIT LINES: Connects data sequentially to reveal the phase trajectory loops
            \addplot3[
                mesh,
                line width=0.8pt,
                mark=none,
                scatter,
                scatter src=explicit,
                shader=interp
            ] table[
                col sep=comma,
                x=eta_1,
                y=eta_2,
                z=eta_3,
                meta=attractor_spin
            ] {../../data/outputs/regime_scatterplots_all.csv};

        \end{axis}
    \end{tikzpicture}
    \caption{3D attractor state orbit trajectory. Color variation tracks Attractor Spin ($\Omega$), separating shear breakout acceleration (cool blue) from radiative collapse configurations (warm red).}
    \label{fig:attractor_ALL}
\end{figure}

\subsection{Physical Coordinate Interpretation}

This phase-space view acts as the compressed state engine of the atmospheric boundary layer (ABL), projecting high-dimensional tower profiles into three dominant, orthogonal coordinates:

\begin{itemize}
    \item $\eta_1$ (\textbf{Bulk Background Mode}): Captures the mean inversion strength, nocturnal cooling depth, and overall boundary-layer depth evolution.
    \item $\eta_2$ (\textbf{Shear \& LLJ Mode}): Captures shear production, alongside low-level jet (LLJ) speed and height modulation.
    \item $\eta_3$ (\textbf{Vertical Structural Curvature Mode}): Captures submesoscale gravity-wave excitation and the precursor structural signatures of intermittent turbulence events.
\end{itemize}

\subsubsection{Geometric Reading Guide}
\begin{itemize}
    \item \textbf{Tight closed loops} indicate near-repeatable diurnal cycling (convective daytime mixing to stable nighttime recovery and back).
    \item \textbf{Broad, tangled trajectories} indicate non-stationary forcing, intermittency, or mesoscale disruption.
    \item \textbf{Color evolution by $\Omega(t)$} isolates transition directionality: positive spin values (warm tones) indicate stable collapse paths; negative values (cool tones) mark turbulent regeneration breakouts.
\end{itemize}

\subsection{Trajectory Metrics}

The mean low-rank coefficients across all samples are:
\begin{align}
    \langle \eta_1 \rangle &= 0.574 \\
    \langle \eta_2 \rangle &= 0.284 \\
    \langle \eta_3 \rangle &= -0.077
\end{align}

The magnitude of the attractor vector $\|\boldsymbol{\eta}\|$ evolves continuously throughout the campaign, reflecting the transient structural adjustments driven by diurnal heating and nocturnal radiative cooling cycles.

Operationally, peaks in $\|\boldsymbol{\eta}\|$ often coincide with rapid boundary-layer reconfiguration, while spin excursions identify windows where similarity-based single-regime assumptions are least reliable.

\subsection{Information-Theoretic Manifold Complexity Evaluation}

To preserve algorithmic invariance across divergent energy scales, rank truncation is enforced with a machine-precision floor criterion ($s_i \ge \epsilon_{\text{mach}} s_1$). This execution yields:

\begin{itemize}
    \item \textbf{Constrained modes:} $N_c = 3$
    \item \textbf{Unconstrained nullspace modes:} $N_0 = 0$
    \item \textbf{Condition number:} $\kappa = 6.85$
    \item \textbf{Shannon effective dimension:} $D_{\text{eff}} = 2.0537$
\end{itemize}

The entropy-derived $D_{\text{eff}} = 2.0537$ quantifies how many modes are actively carrying structure in practice, rather than only counting nominal basis size. Larger values indicate broader modal participation and lower compressibility of the observed boundary-layer state.

\label{sec:attractor_ALL}

\newpage

% AUTO-GENERATED Regime Structural Analysis
% Rendered by Mustache template from curated diagnostics
% DO NOT EDIT MANUALLY

\section{Spectral Geometry and Regime Structural Analysis}

\subsection{Mathematical Foundation}

The decomposition of the boundary layer into regimes relies on the spectral properties of the low-rank attractor subspace. We define three primary diagnostic metrics:

\begin{enumerate}
    \item \textbf{Singular Value Entropy} $H = -\sum_i p_i \log(p_i)$ where $p_i = \sigma_i^2 / \sum_j \sigma_j^2$
    \item \textbf{Spectral Curvature} $\chi_N = \frac{\eta_3}{\eta_1 + \eta_2 + \epsilon}$ (normalized vertical component)
    \item \textbf{Attractor Magnitude} $\|\boldsymbol{\eta}\| = \sqrt{\eta_1^2 + \eta_2^2 + \eta_3^2}$
\end{enumerate}

\subsection{Campaign Context}

\begin{itemize}
    \item \textbf{Campaign:} ALL
    \item \textbf{Baseline:} v0.1 (spectralbl-attractor)
    \item \textbf{Temporal Coverage:} 59149.5 hours
    \item \textbf{Sample Count:} 6682 observations
\end{itemize}

\subsection{Entropy Evolution}

The mean singular value entropy for this campaign is:
\[
    H_{\text{mean}} = 0.223
\]

This value reflects the typical balance between coherent, low-dimensional structure and multi-scale turbulent activity. Nocturnal periods show reduced entropy (stratified regime), while daytime transitions display elevated entropy (mixing regime).

\subsection{Structural Independence Verification}

To verify that the low-rank decomposition captures physically independent aspects of boundary-layer dynamics, we contrast spectral geometry against classical stability diagnostics.

The gradient Richardson number can be interpreted as
\[
\mathrm{Ri}_g \sim \frac{\text{buoyant suppression}}{\text{shear production}}.
\]
In practice, low $\mathrm{Ri}_g$ indicates active shear-driven turbulence, while large $\mathrm{Ri}_g$ indicates strongly stable conditions where continuous turbulence is difficult to sustain.

However, in very stable boundary layers, $\mathrm{Ri}_g$ alone is often too blunt: it indicates suppression, but not whether residual energy appears as intermittent turbulence, submesoscale motions, or wave-dominated structure.

The spectral curvature metric $\chi_N$ addresses this by tracking geometric structure in the reduced attractor space rather than relying only on local gradients.

When $\mathrm{Ri}_g$ exceeds the critical threshold $(0.2-0.250)$, classical surface-layer similarity theory predicts that continuous turbulence should collapse entirely. In practice, $\mathrm{Ri}_g$ then saturates and becomes uninformative at precisely the moment when the boundary layer is most structurally active---dominated by intermittent bursting, gravity-wave excitation, and slow turbulent kinetic energy decay that gradient metrics cannot resolve.

The scatter plot shown in Figure~\ref{fig:spectral_orthogonality_all} demonstrates this directly. While $\mathrm{Ri}_g$ stalls at high values, the vertical coordinate $\eta_3$ remains dynamically resolved throughout the same stability range, mapping the active wave field and intermittent bursting events that $\mathrm{Ri}_g$ misses entirely.

Scientific interpretation of orthogonality:
\begin{itemize}
    \item If $\chi_N$ were redundant with $\mathrm{Ri}_g$, points would collapse toward a single diagonal trend.
    \item Distinct cross-patterns or separated blocks indicate independent information content: the spectral decomposition is resolving physics that local gradient scalings cannot.
    \item This orthogonality is the core scientific justification for the SpectralBL framework: it provides structural regime information in conditions where MOST-based parameterizations break down.
\end{itemize}

\begin{figure}[htbp]
    \centering
    \begin{tikzpicture}
        \begin{axis}[
            width=0.86\textwidth,
            height=0.48\textwidth,
            grid=major,
            xlabel={$\eta_1$},
            ylabel={$\eta_2$},
            title={Spectral Orthogonality Projection (ALL)},
            colormap/plasma,
            colorbar,
            point meta=explicit,
            colorbar style={title={$\eta_3$}},
            mark size=1.8pt,
            tick label style={font=\small},
            label style={font=\small},
            title style={font=\small}
        ]
            \addplot[
                only marks,
                scatter,
                scatter src=explicit
            ] table[
                col sep=comma,
                x=eta_1,
                y=eta_2,
                meta=eta_3
            ] {../../data/outputs/regime_scatterplots_all.csv};
        \end{axis}
    \end{tikzpicture}
    \caption{Projected attractor coordinates with orthogonality structure visualized via $\eta_3$ color encoding.}
    \label{fig:spectral_orthogonality_all}
\end{figure}

\subsection{Non-Orthogonal Scale Coupling Accounts}

Because smooth partition-of-unity masks ($\psi_M,\psi_W,\psi_T$) retain localized overlap to stabilize regime transitions, total energy tracking is not strictly additive at every instant.

For this execution, the off-diagonal interaction proxy is
\[
\mathcal{I}_{ij} = 0.000000.
\]

This term is tracked to quantify cross-regime leakage (wave-turbulence exchange) during active stratification shifts and to preserve auditability of coupling behavior in the reduced-order representation.

\paragraph{Assumption Note}
Here $\mathcal{I}_{ij}$ is a variability proxy derived from entropy dynamics, not a closed-form energetic decomposition term; it is intended for comparative run diagnostics rather than absolute budget closure.

\label{sec:regime_all}


\newpage

% AUTO-GENERATED Pipeline Verification and Analytical Conclusions
% Rendered by Mustache template from curated diagnostics
% DO NOT EDIT MANUALLY

\section{System Diagnostics and Definitive Conclusions}

\subsection{Pipeline Execution Verification}

The execution of the \texttt{SpectralBL-Analytics} pipeline across divergent data scales demonstrates the mathematical robustness of the low-rank attractor manifold projection. The framework successfully achieves structural decomposition without reliance on local gradient Richardson number ($\text{Ri}_g$) formulations, which typically fail under stable conditions.

\subsection{Comparative Information-Theoretic Summary}

Based on the processed analytics run, the structural characteristics of the underlying dataset are summarized below:

\begin{table}[htbp]
    \centering
    \small
    \begin{tabular}{l p{9cm}}
        \hline
        	extbf{Diagnostic Metric} & \textbf{Pipeline Evaluation Summary for \texttt{ CASES-99 }} \\ \hline
        Data Footprint & Analyzed 6682 observations across a temporal scale of 59149.5 hours. \\
        Manifold Complexity & Mean singular value entropy registers at $H_{\text{mean}} = 0.223$, yielding a Shannon effective dimension of $D_{\text{eff}} = 2.0537$. \\
        Numerical Stability & The condition number $\kappa = 6.85$ confirms a well-conditioned low-rank basis projection ($N_0 = 0$ nullspace modes). \\ \hline
    \end{tabular}
    \caption{Key pipeline execution and manifold metrics for the current campaign.}
    \label{tab:pipeline_summary_CASES-99}
\end{table}

\subsection{Definitive Scientific Conclusions}

By mapping the boundary layer dynamics into the orthogonal coordinates $(\eta_1, \eta_2, \eta_3)$, this analysis establishes the following definitive conclusions based on the campaign-specific footprint:


\begin{itemize}
    \item \textbf{Attractor Collapse Identification:} The low mean entropy and compressed effective dimension provide empirical evidence of attractor collapse behavior. Across its long timeline, \texttt{CASES-99} frequently drops into highly ordered, low-dimensional states.
    \item \textbf{Dominance of Non-Local Regimes:} The structural decomposition indicates persistent intermittent shear-burst behavior, explaining why strictly local similarity theories degrade in this regime.
    \item \textbf{Validation of HOST-Like Transition Behavior:} The coordinate tracking resolves non-linear transitions where turbulence ceases to be continuously gradient-driven. While local metrics such as $\text{Ri}_g$ saturate, spectral geometry remains dynamically informative.
\end{itemize}

\subsection{Operational Framework Recommendation}

These signatures confirm that \texttt{GABLS3} and \texttt{CASES-99} should be used as complementary axes in the verification pipeline:
\begin{equation}
    \mathcal{M}_{\text{validation}} = \mathbf{f}\big(\underbrace{\text{GABLS3}}_{\text{Trajectory Fidelity (Process)}},\ \underbrace{\text{CASES-99}}_{\text{Ergodic Sampling (Climatology)}}\big)
\end{equation}

While \texttt{GABLS3} evaluates transient diurnal arc fidelity, \texttt{CASES-99} provides statistical bounds for model behavior in deep stratification, IGW coupling, and collapse-prone states over climatological windows.

\label{sec:conclusions_CASES-99}


\newpage

% templates/bifurcation_spectrum.tex.mustache
% AUTO-GENERATED Stage 5 bifurcation panel
% Rendered by scripts/build_campaign_report.jl
% DO NOT EDIT MANUALLY

\section{Physical Interpretation of the Continuation Parameter}

The continuation parameter $\gamma$ represents the effective atmosphere--surface coupling strength inferred from the sparse structural operator identified in Stage~4.
Small values of $\gamma$ correspond to weak turbulent transport and strongly decoupled nocturnal conditions, where the surface inversion evolves in relative isolation from the overlying residual layer.
Larger values of $\gamma$ represent enhanced vertical exchange between the surface inversion and the residual layer above, sustaining coherent turbulent structures through the column.

Parameter continuation therefore traces the geometric deformation of the reconstructed attractor manifold as the atmospheric state transitions between coupled and decoupled regimes.
The leading eigenvalue $\alpha(\gamma) = \max\,\Re(\lambda_i)$ serves as a scalar indicator of manifold stability: negative values indicate a contracting, stable attractor; the zero crossing at $\gamma_c$ marks the onset of structural instability; and the imaginary part $\beta(\gamma) = |\mathrm{Im}(\lambda_\mathrm{dom})|$ tracks the intrinsic oscillatory frequency of the emerging mode.

\section{Stage 5 Bifurcation Spectrum}


\noindent\textit{Stage 5 branch artifacts were not available for this campaign at report render time. Expected source: \path{../../data/outputs/stage5_bifurcation_branches_all.csv}.}



\noindent\textit{Cross-campaign Stage 5 comparison table omitted because one or more required summary manifests were not available at render time.}


\newpage

\appendix

\section{Pipeline Metadata}

\begin{itemize}
    \item Source code: \url{https://github.com/davidengland/SpectralBL-Analytics}
    \item Build timestamp: \today{}
    \item Supported campaigns: CASES-99, GABLS3, AMERIFLUX (60+ sites), SMEAR, NEON, ICOS, SHEBA
    \item Mathematical basis: p-FEM spectral projection and scale-invariant rank truncation
    \item Templating: Mustache -- active
    \item Figure cache: TikZ external compilation -- \texttt{tikz-cache/}
\end{itemize}

\end{document}
